\documentclass[12pt,a4paper]{article}
\usepackage{amsmath,amssymb}
\usepackage{booktabs}
\usepackage{geometry}
\usepackage{graphicx}
\usepackage{float}
\usepackage{listings}
\usepackage{xcolor}
\usepackage{algorithm}
\usepackage{algpseudocode}
\usepackage{titlesec}
\usepackage{setspace}
\usepackage{ragged2e}

%----------------------------------------
% Page Setup for A4
%----------------------------------------
\geometry{
  a4paper,
  top=1in,
  bottom=1in,
  left=1in,
  right=1in
}

\setstretch{1.2}

%----------------------------------------
% Section title formatting
%----------------------------------------
\titleformat{\section}
  {\normalfont\bfseries\fontsize{14}{14}\selectfont}
  {\thesection}{1em}{}

%----------------------------------------
% Code listing style
%----------------------------------------
\lstdefinestyle{matlabstyle}{
  language=Matlab,
  basicstyle=\ttfamily\footnotesize,
  keywordstyle=\color{blue}\bfseries,
  commentstyle=\color{green!60!black},
  stringstyle=\color{orange!80!black},
  frame=none,
  numbers=none,
  breaklines=true,
  showstringspaces=false,
  tabsize=2
}

%----------------------------------------
\begin{document}

%----------------------------------------
% Cover Page
%----------------------------------------
\begin{titlepage}
\centering

\vspace*{0.5cm}
\includegraphics[width=0.22\textwidth]{diu_logo.png}\par
\vspace{0.5cm}

{\Large \textbf{Department of Computer Science and Engineering} \par}
\vspace{0.2cm}
{\Large \textbf{Dhaka International University} \par}

\vspace{0.8cm}

{\normalsize
Semester: 12th \par
Lab Report No: 02 \par
Course Code: CSE-420 \par
Course Title: Image Processing Lab \par
Report Title: Image Noise Addition Using imnoise and Manual Methods in MATLAB \par
}

\vspace{1cm}

{\Large \textbf{Submitted By} \par}
\vspace{0.3cm}

{\normalsize
Rasheduzzaman Rakib \par
Reg. No.: CS-D-77-22-120080 \par
Roll: 37 \par
Batch: D-77 \par
Department of Computer Science and Engineering \par
Dhaka International University \par
}

\vspace{1cm}

{\Large \textbf{Submitted To} \par}
\vspace{0.3cm}

{\normalsize
Md. Muksit Ul Islam \par
Assistant Professor \par
Department of Computer Science and Engineering \par
Dhaka International University \par
}

\vfill

{\normalsize Submission Date: 25th September 2026 \par}
{\normalsize Summer 2026 \par}

\end{titlepage}


%----------------------------------------
% Lab Report
%----------------------------------------

\section{Name}

\textbf{Image Noise Addition Using imnoise and Manual Methods in MATLAB}\\[0.3em]
Topics covered: Gaussian Noise, Salt \& Pepper Noise, Poisson Noise,
Manual Rayleigh Noise, Manual Gamma Noise, Manual Exponential Noise,
Manual Uniform Noise.

%----------------------------------------

\section{Description}

This lab demonstrates the addition of seven different types of noise to a
grayscale image in MATLAB. The built-in image peppers.png is first converted
to grayscale, then each noise type is applied and displayed in a subplot grid.
Three noises are added using MATLAB's imnoise function; the remaining four
are constructed manually using probability distributions.

\begin{enumerate}
  \item \textbf{Gaussian Noise:} Added via imnoise with the 'gaussian' model.
        It follows a normal distribution and manifests as a uniform graininess spread across the entire image.

  \item \textbf{Salt \& Pepper Noise:} Added via imnoise with density $d = 0.05$.
        Randomly sets 5\% of pixels to either 0 (pepper) or 255 (salt),
        creating scattered black and white dots.

  \item \textbf{Poisson Noise:} Added via imnoise with the 'poisson' model.
        Noise variance equals the local pixel intensity, making it
        signal-dependent (brighter areas have more variance).

  \item \textbf{Manual Rayleigh Noise:} Generated using the inverse CDF method.
        Scaled by 20 and added to the grayscale image.

  \item \textbf{Manual Gamma Noise:} Drawn from a Gamma distribution using
        gamrnd(2, 1), scaled by 20, and added to the pixel intensities.
        Produces asymmetric, positively-skewed noise.

  \item \textbf{Manual Exponential Noise:} Drawn from an exponential distribution
        using exprnd(1), scaled by 20, and added to the image.
        Results in heavy right-tailed intensity shifts.

  \item \textbf{Manual Uniform Noise:} Drawn from $\mathcal{U}(0,1)$ using rand,
        scaled by 20, and added uniformly across all pixels, causing a
        mild, even brightness increase with random variation.
\end{enumerate}

%----------------------------------------

\section{Code}

\begin{lstlisting}[style=matlabstyle]
I = imread('peppers.png');
gr = rgb2gray(I);

% 1. gaussian noise
gaussian = imnoise(gr, 'gaussian');
subplot(3,3,1);
imshow(gaussian);
title('gaussian noise');

% 2. salt and pepper noise
sp = imnoise(gr, 'salt & pepper', 0.05);
subplot(3,3,2);
imshow(sp);
title('salt and pepper noise');

% 3. poisson noise
poisson = imnoise(gr, 'poisson');
subplot(3,3,3);
imshow(poisson);
title('poisson noise');

% 4. manual rayleigh noise
rayleigh = rand(size(gr));
rayleigh = sqrt(-2 * log(1 - rayleigh));
rayleigh = uint8(double(gr) + rayleigh * 20);
subplot(3,3,4);
imshow(rayleigh);
title('rayleigh noise');

% 5. manual gamma noise
gamma = gamrnd(2, 1, size(gr));
gamma = uint8(double(gr) + gamma * 20);
subplot(3,3,5);
imshow(gamma);
title('gamma noise');

% 6. manual exponential noise
exponential = exprnd(1, size(gr));
exponential = uint8(double(gr) + exponential * 20);
subplot(3,3,6);
imshow(exponential);
title('exponential noise');

% 7. manual uniform noise
uniform = rand(size(gr));
uniform = uint8(double(gr) + uniform * 20);
subplot(3,3,7);
imshow(uniform);
title('uniform noise');
\end{lstlisting}

%----------------------------------------

\section{Output}

Figure~\ref{fig:output} shows the seven noise types applied to the grayscale
peppers.png image.

\begin{figure}[H]
  \centering
  \includegraphics[width=0.82\textwidth]{Figure_1.png}
  \caption{Seven noise models applied to grayscale peppers.png test image.}
  \label{fig:output}
\end{figure}

Figure~\ref{fig:output} confirms the distinct visual characteristics of each
noise model: Gaussian (1) produces uniform grain; Salt \& Pepper (2) shows
isolated white/black dots; Poisson (3) is visually subtle as it is
signal-dependent; the manually generated noises (4--7) produce a brightening
effect due to positive-valued additive offsets scaled by 20.

\end{document}
